Our fixed-reference audit turns generative diversity assessment into a framework for locating, quantifying, and repairing distributional gaps. By separating regional probability, conditional spread, and reference support, it makes changes in a generator’s output distribution interpretable and actionable. Finite-sample certificates establish the magnitude of regional mass changes, explicit sample requirements define the audit’s resolution, and controlled interventions test whether a diagnosis leads to improvement on independent outcomes. The experiments demonstrate this connection at matched computational budgets. Covariance initialization reduces two-call CIFAR-10 KID by 32.0\% and improves real-data coverage without additional denoiser calls, with gains retained against CIFAR-10.1. Ordinary continuation reduces MNIST class-mass error by 58.7\% at equal training updates and improves total variation across all classifier outcomes, including rejection. Gaussian analysis separates initialization from integration error, while crossed probes distinguish sampler-sensitive deficits from persistent depletion. Across four text-to-image families, the audit reveals diversity reversals caused by pooling, representation, and teacher choice. Together, these results establish a practical link between statistically grounded distributional diagnosis and targeted generator improvement, giving practitioners a concrete basis for deciding what to change.